\documentclass[aspectratio=169]{beamer}

\usepackage{fontspec}
\usepackage{polyglossia}
\usepackage{color}
\usepackage{listings}
\setdefaultlanguage{russian}
\usepackage{csquotes}
\usetheme{metropolis}
\metroset{background=light}
\metroset{sectionpage=none}
\metroset{numbering=fraction}

\definecolor{ssaublue}{RGB}{77,144,205}
\setbeamercolor{frametitle}{fg=white, bg=ssaublue}
\setbeamercolor{title separator}{fg=ssaublue}

\setmainfont[
    Path = ../fonts/elektratextpro/,
    Extension = .otf,
    UprightFont = elektratextpro,
    BoldFont = elektratextpro_bold,
    ItalicFont = elektratextpro_italic,
    BoldItalicFont = elektratextpro_bolditalic
]{elektratextpro}

\setsansfont[
    Path = ../fonts/elektratextpro/,
    Extension = .otf,
    UprightFont = elektratextpro,
    BoldFont = elektratextpro_bold,
    ItalicFont = elektratextpro_italic,
    BoldItalicFont = elektratextpro_bolditalic
]{elektratextpro}

\setmonofont[
    Path = ../fonts/elektratextpro/,
    Extension = .otf,
    UprightFont = elektratextpro,
    BoldFont = elektratextpro_bold,
    ItalicFont = elektratextpro_italic,
    BoldItalicFont = elektratextpro_bolditalic
]{elektratextpro}

\newfontfamily\cyrillicfont[
    Path = ../fonts/elektratextpro/,
    Extension = .otf,
    UprightFont = elektratextpro,
    BoldFont = elektratextpro_bold,
    ItalicFont = elektratextpro_italic,
    BoldItalicFont = elektratextpro_bolditalic
]{elektratextpro}

\newfontfamily\cyrillicfontsf[
    Path = ../fonts/elektratextpro/,
    Extension = .otf,
    UprightFont = elektratextpro,
    BoldFont = elektratextpro_bold,
    ItalicFont = elektratextpro_italic,
    BoldItalicFont = elektratextpro_bolditalic
]{elektratextpro}

\newfontfamily\cyrillicfonttt[
    Path = ../fonts/elektratextpro/,
    Extension = .otf,
    UprightFont = elektratextpro,
    BoldFont = elektratextpro_bold,
    ItalicFont = elektratextpro_italic,
    BoldItalicFont = elektratextpro_bolditalic
]{elektratextpro}

% Настройка листингов для Beamer
\lstset{
    basicstyle=\ttfamily\footnotesize,
    commentstyle=\color{gray},
    columns=fullflexible,
    keepspaces=true,
    breaklines=false,
    extendedchars=true,
    texcl=true,
    showstringspaces=false
}

\title{Git в DevOps}
\date{17 сентября 2025 г.}

\begin{document}
    \maketitle

    \begin{frame}{Эволюция роли Git}
        \begin{columns}
            \begin{column}{0.48\textwidth}
                \textbf{Git ранее}
                \begin{itemize}
                    \item Хранилище исходного кода
                    \item Версионирование файлов
                    \item Работа с ветками
                    \item Совместная разработка
                \end{itemize}
            \end{column}

            \begin{column}{0.48\textwidth}
                \textbf{Git в DevOps}
                \begin{itemize}
                    \item Центр управления системой
                    \item Source of Truth для системы
                    \item Триггер автоматизации
                    \item Аудит изменений
                \end{itemize}
            \end{column}
        \end{columns}

        \vspace{0.5cm}

        \textbf{Git репозиторий} становится местом хранения состояния \textbf{всей системы}
    \end{frame}

    \begin{frame}{Что хранится в DevOps репозитории?}
        \begin{center}
            \large
            \texttt{my-project/}
        \end{center}

        \begin{itemize}
            \item \texttt{src/} \textcolor{gray}{— исходный код приложения}
            \item \texttt{config/} \textcolor{gray}{— файлы конфигурации}
            \item \texttt{scripts/} \textcolor{gray}{— скрипты развертывания}
            \item \texttt{docs/} \textcolor{gray}{— документация и инструкции}
            \item \texttt{tests/} \textcolor{gray}{— тесты и проверки}
        \end{itemize}

        \vspace{0.2cm}

        \textbf{Принцип:} если это влияет на систему — это должно быть в Git

        \footnotesize{\textcolor{gray}{\url{https://octopus.com/blog/what-is-everything-as-code}}}
    \end{frame}

    \begin{frame}{Демонстрация: проект в DevOps стиле}

        \textbf{Что покажем:}
        \begin{itemize}
            \item Создание репозитория с веб-сайтом
            \item Конфигурация nginx в Git
            \item Скрипт автоматического развертывания
            \item Изменение сайта через Git
        \end{itemize}
        
        \vspace{0.2cm}
        
        \textbf{Демо-репозиторий:} \\
        \href{https://github.com/prafdin/devops-demo-website}{\texttt{https://github.com/prafdin/devops-demo-website}}
    \end{frame}

    \begin{frame}{Git как триггер автоматизации}
        \begin{center}
            \large
            \textbf{События в Git запускают процессы}
        \end{center}
        
        \vspace{0.5cm}
        
        \begin{columns}
            \begin{column}{0.48\textwidth}
                \textbf{Традиционный подход:}
                \begin{itemize}
                    \item Ручное развертывание
                    \item Разделенные процессы
                    \item Человеческий фактор
                    \item Нет связи код <-> система
                \end{itemize}
            \end{column}
            
            \begin{column}{0.48\textwidth}
                \textbf{DevOps подход:}
                \begin{itemize}
                    \item git push → автодеплой
                    \item git tag → релиз
                    \item Pull Request → тестирование
                    \item Коммит = событие системы
                \end{itemize}
            \end{column}
        \end{columns}
        
        \vspace{0.5cm}
        
        \textbf{Принцип:} каждое изменение в Git влияет на всю систему
    \end{frame}

    \begin{frame}{Демонстрация: GitHub Webhook}
        \textbf{Что покажем:}
        \begin{itemize}
            \item Настройка GitHub webhook
            \item Настройка доступа к ВМ через реверс прокси
            \item Автоматический запуск деплоя по событию в репозитории
        \end{itemize}
        
        \vspace{0.2cm}
        \textbf{Демо-репозиторий:} \\
        \href{https://github.com/prafdin/devops-demo-website/tree/webhook-demo}{\texttt{https://github.com/prafdin/devops-demo-website/tree/webhook-demo}}
    \end{frame}

    \begin{frame}{Git ветки = Окружения системы}
        \begin{center}
            \large
            \textbf{Каждому окружению - своя ветка}
        \end{center}
        
        \vspace{0.5cm}
        
        \begin{itemize}
            \item \texttt{main} → Продакшн окружение
            \item \texttt{develop} → Тестовое окружение
            \item \texttt{feature/*} → Локальная разработка / персональное тестовое окружение
        \end{itemize}
        
        \vspace{0.2cm}

        Ветки в git - одно из средств управлениями окружениями в DevOps репозитории
    \end{frame}

    \begin{frame}{Демонстрация: Environments через Git}
        \textbf{Что покажем:}
        \begin{itemize}
            \item Webhook сервер определяет окружение по ветке из payload
            \item Деплой скрипт умеет работать с разными окружениями
            \item Изменения в develop отображаются в app-test.prafdin.course.prafdin.ru, а изменения в основное ветке в app.prafdin.course.prafdin.ru
        \end{itemize}
        
        \vspace{0.2cm}
        \textbf{Демо-репозиторий:} \\
        \href{https://github.com/prafdin/devops-demo-website/tree/environments-demo}{\texttt{https://github.com/prafdin/devops-demo-website/tree/environments-demo}}
    \end{frame}

    \begin{frame}{Проблемы текущего подхода}
        \begin{itemize}
            \item Webhook сервер нужно запускать и обновлять вручную.
            \item Под Webhook сервер нужна отдельная ВМ.
            \item Nginx устанавливается вручную
            \item Отсутствие наглядности: логи доступны только на ВМ, в репозитории не видно статуса выполнения.
        \end{itemize}
    \end{frame}

    \begin{frame}{Что запомнить}
        \begin{itemize}
            \item Git центральный элемент в DevOps инструментах
            \item События в Git - триггер автоматизации
            \item В Git находится не только код приложения, но и все что влияет на систему
        \end{itemize}

        \vspace{1cm}

        \begin{center}
            \textbf{Вопросы?}
        \end{center}
    \end{frame}
\end{document}